% Introduction draft (session 2); merged into main.tex.
\section*{Introduction}

A robot that grasps a known object usually first estimates the object's six-dimensional pose and then plans the grasp on the object's model~\cite{xiang2018posecnn,tremblay2018dope,deng2020selfsup}. Pose estimators are developed and ranked on benchmarks such as LINEMOD~\cite{hinterstoisser2012} and BOP~\cite{hodan2018bop} by geometric criteria: the average distance between model points under the estimated and the true pose (ADD, or ADD-S for symmetric objects) must fall below a tenth of the object diameter. On LINEMOD, current methods pass this test on nearly every image; FoundationPose, for example, passes it on \FpAdd\ of the test scenes used here~\cite{wen2024foundationpose}. Whether a grasp planned on such a pose actually lifts the object is not measured by these benchmarks, and the threshold was not designed to answer it~\cite{hodan2016eval}: a grasp can fail under a pose error the metric accepts, and it can fail for reasons no pose contains, such as friction or the position of the centre of mass.

Grasp-learning methods take the opposite route and predict grasp success directly from images or point clouds~\cite{mahler2017dexnet,morrison2018ggcnn,tenpas2017gpd,sundermeyer2021contact,fang2020graspnet,fang2023anygrasp}, and reviews of the field~\cite{bohg2014survey,newbury2023review} describe a steady move toward such outcome-supervised models. These methods are evaluated on their own data sets, with their own grippers and success criteria, so the two traditions are rarely compared on the same images, the same candidate grasps and the same outcomes.

We make this comparison possible on real images. LM-Grasp pairs every real LINEMOD test image of the 13 standard objects (\NScenes\ RGB-D scenes) with 16 candidate grasps proposed from the observation and with outcomes obtained by executing each candidate on the object's scanned mesh at its annotated pose in a physics simulator, with friction, mass and centre of mass drawn per scene and never shown to any method (Fig.~\ref{fig:overview}). The images are real; the grasp outcomes are simulated. On this benchmark we evaluate two routes from the same observation to a grasp decision: pose-then-plan pipelines that give a pose to a planner holding the object model and the simulator, and manipulation-sufficient perception, which learns from the image a compact code from which the success of any candidate grasp can be read off~\cite{tishby2000ib,alemi2017vib,pacelli2020}. Because the code is formed before a grasp is chosen, one code scores every candidate, and its size is the payload a robot would send to, or receive from, an edge server~\cite{kehoe2015cloud,shao2022taskcomm}.

Our contributions are fourfold. (i) We release LM-Grasp, an outcome-level protocol on real LINEMOD images with \NRollouts\ simulated grasp executions. (ii) We show that the ADD(-S) criterion is a weak proxy for grasp success: pose errors that it accepts on every scene still cost up to \NoiseSatDrop\ in top-1 grasp success, and FoundationPose, at \FpAdd\ pose accuracy, plans grasps that succeed \FpPlanTop\ of the time against \GtPlanTop\ with the exact pose. (iii) We show that perception trained on grasp outcomes, with no object model or simulator at decision time, reaches \EnsTopMean\ top-1 success and selects grasps that succeed \EnsSelQMean\ of the time on the quarter of scenes it is most confident about, and that its decision-relevant content fits into a few tens of bits. (iv) We compare a stochastic, rate-regularised code (MSP) with a deterministic code compressed after training, a deep ensemble, a grasp-patch network with the Dex-Net 2.0 architecture and learned and model-based pose pipelines under matched data, candidates and training budget, with five seeds and paired tests, and we report where each one stands.
